%La línea de abajo es para quitar encabezado
%\thispagestyle{plain}

\chapter*{Introducción}
\markboth{Introducción}{Introducción}
\addcontentsline{toc}{chapter}{Introducción}

En los últimos años, especialmente con la masificación de las redes sociales y plataformas digitales, los procesos electorales han dejado de analizarse únicamente a través de encuestas tradicionales para incorporar nuevas fuentes de información que reflejan la opinión pública en tiempo real. Periódicos digitales, plataformas audiovisuales como YouTube y redes sociales como Twitter/X generan grandes volúmenes de datos que capturan tanto el discurso mediático como la percepción ciudadana. Entre estos entornos digitales, han surgido nuevas dinámicas de comunicación política, así como nuevas oportunidades de investigación orientadas a comprender cómo se construye y evoluciona la opinión pública durante campañas electorales.

A diferencia de estudios previos, en donde gran parte de los enfoques se centra en el análisis de una sola fuente de datos o en modelos tradicionales de procesamiento de lenguaje natural, la principal motivación de esta investigación es proponer un enfoque integral de Inteligencia Electoral que permita analizar de manera conjunta múltiples fuentes heterogéneas. Bajo este escenario, se busca superar las limitaciones asociadas al uso de una única plataforma —como los sesgos de representación o la falta de contexto— mediante la integración de datos provenientes de prensa digital, contenido audiovisual y redes sociales. El aporte principal consiste en el diseño de un sistema basado en Modelos de Lenguaje de Gran Escala (LLMs) que no solo realiza minería de sentimientos, sino que también detecta dinámicamente temas relevantes y modela las relaciones entre actores, fuentes y tópicos a través de un grafo de conocimiento temporal.

Asimismo, la propuesta incorpora un enfoque innovador al considerar explícitamente la diferencia entre el discurso mediático y la reacción ciudadana, permitiendo analizar la brecha entre ambos como un indicador clave en el comportamiento electoral. Esta aproximación se complementa con el uso de técnicas avanzadas como clustering inductivo de temas y sistemas GraphRAG, los cuales permiten realizar consultas en lenguaje natural sobre la información estructurada y no estructurada, facilitando la generación de insights complejos y contextualizados.

En el Capítulo I, se describe la realidad problemática de la investigación en el contexto de los procesos electorales y los ecosistemas digitales. Asimismo, se formulan los problemas, objetivos, hipótesis, justificación y delimitación del estudio.

En el Capítulo II, se detallan los antecedentes principales relacionados con el análisis de sentimiento político, el uso de modelos de lenguaje, el análisis multi-fuente y los grafos de conocimiento en contextos electorales. Posteriormente, se presenta el marco teórico que sustenta la investigación, abordando conceptos de procesamiento de lenguaje natural, inteligencia artificial, análisis de opinión y comunicación política digital.

En el Capítulo III, se describe el diseño metodológico de la investigación, incluyendo el tipo, enfoque, población y muestra. Asimismo, se detalla la arquitectura del sistema propuesto, basada en un enfoque tipo Medallion, así como el pipeline de procesamiento de datos, desde la ingesta hasta la construcción del grafo de conocimiento. Finalmente, se expone la metodología de evaluación y el cronograma de actividades.

En el Capítulo IV, se desarrolla la implementación de la solución propuesta, incluyendo la integración de fuentes de datos, el uso de LLMs para el análisis semántico, la detección de temas mediante clustering y la construcción del grafo de conocimiento, así como el sistema de consulta basado en GraphRAG.

En el Capítulo V, se presentan y analizan los resultados obtenidos, evaluando el desempeño del sistema en términos de detección de sentimientos, coherencia temática y capacidad de generación de insights. Asimismo, se discuten los hallazgos relacionados con la dinámica de la opinión pública y la brecha entre medios y ciudadanía.

En el Capítulo VI, se exponen las conclusiones de la investigación y se plantean recomendaciones para futuras líneas de trabajo, incluyendo posibles mejoras en la arquitectura del sistema y la ampliación del análisis a otros contextos electorales.

La investigación concluye con las referencias bibliográficas utilizadas y los anexos que complementan la información presentada a lo largo del documento.